% ---- pg-406 （书页 394） ----

\begin{mvtable}
\begin{tabular}{lccc}
\toprule
\multirow{2}{*}{变\hspace{1.5em}量} & \multicolumn{2}{c}{旋转因子载荷估计} & \multirow{2}{*}{共性方差 $\hat{h}_i^{*2}=\hat{h}_i^2$}\\
\cmidrule(lr){2-3}
& $F_1^*$ & $F_2^*$ & \\
\midrule
1.\ 盖尔语 & 0.369 & \textbf{0.594} & 0.490\\
2.\ 英语 & 0.433 & \textbf{0.467} & 0.406\\
3.\ 历史 & 0.211 & \textbf{0.558} & 0.356\\
4.\ 算术 & \textbf{0.789} & 0.001 & 0.623\\
5.\ 代数 & \textbf{0.752} & 0.054 & 0.569\\
6.\ 几何 & \textbf{0.604} & 0.083 & 0.372\\
\bottomrule
\end{tabular}
\phantomsection\label{tab:9.6}\mvtabcap{表 9.6}
\end{mvtable}

由于 $\hat{\mathbf{L}}\hat{\mathbf{L}}'=\hat{\mathbf{L}}\mathbf{T}\mathbf{T}'\hat{\mathbf{L}}'=\hat{\mathbf{L}}^*\hat{\mathbf{L}}^{*'}$, 因此正交旋转不改变估计的共性方差, 并且共性方差是这些矩阵的对角元.

我们指出, \hyperref[fig:9.1]{图 9.1} 表明了坐标的斜交旋转. 一条新轴会通过 $\{1,2,3\}$ 组, 而另一条通过 $\{4,5,6\}$ 组. 斜交旋转的得名是因为它们对应于坐标轴的非刚性旋转, 会产生不相互垂直的新轴. 不过, 很明显, 对于本例斜交因子的阐述与先前给出的对于正交旋转的阐述会是一样的. \hfill$\blacksquare$

凯泽(Kaiser, 见文献[9])曾建议过一个简单结构的分析测度, 称作最大方差准则或正态最大方差准则. 把 $\tilde{l}_{ij}^*=\hat{l}_{ij}^*/\hat{h}_i$ 定义为以共性方差平方根进行标度变换的旋转后系数. 那么此(正态)方差最大过程是选正交变换 $\mathbf{T}$ 使
\phantomsection\label{eq:9-45}
\begin{equation}
V=\frac{1}{p}\sum_{j=1}^{m}\left[\sum_{i=1}^{p}\tilde{l}_{ij}^{*4}-\left(\sum_{i=1}^{p}\tilde{l}_{ij}^{*2}\right)^2\Big/p\right] \tag{9-45}
\end{equation}
尽可能大.

对旋转后系数 $\hat{l}_{ij}^*$ 进行标度变换, 结果使有较小共性方差的变量在决定简单结构时有相对较大的权重. 确定变换 $\mathbf{T}$ 之后, 以 $\hat{h}_i$ 乘载荷 $\tilde{l}_{ij}^*$ 以保存原来的共性方差.

尽管\hyperref[eq:9-45]{式(9-45)}看上去有点可怕, 它却有一个简单的解释. 换言之,
\phantomsection\label{eq:9-46}
\begin{equation}
V\propto\sum_{j=1}^{m}(\text{第 }j\text{ 个因子(标度变换后)载荷平方的方差}) \tag{9-46}
\end{equation}
结果, 最大化 $V$ 对应于把每个因子上的载荷平方尽可能多地“扩展出去”. 因此我们希望在旋转后载荷矩阵 $\hat{\mathbf{L}}^*$ 的任一列上, 找到大的和可以忽略的系数组.

有最大化 $V$ 的算法, 而且最受欢迎的因子分析计算机程序(例如, 统计软件包 SAS, SPSS, BMDP 及 MINITAB)提供最大方差旋转. 正如期望的那样, 用不同解法(主成分、极大似然等)得到的因子载荷的最大方差旋转, 一般说不会是一致的. 如果在旋转中添加公共因子, 那么旋转载荷的模型也可以有明显的改变. 如果一个起控制作用的单因子存在, 用任何正交旋转其含义一般都将被弄模糊. 另一方面, 可将其固定而旋转剩下的因子.

\noindent\textbf{例 9.9}\quad (对消费者偏好数据的旋转载荷)

我们回到例 9.3 中讨论的营销数据. 原始的因子载荷(由主成分法得到)、共性方差及(最大方差)旋转因子载荷, 列在\hyperref[tab:9.7]{表 9.7} 中(见显示板 9.1 中的 SAS 统计软件输出.)

很清楚, 变量 2, 4 和 5 定义因子 1(在因子 1 上有高载荷而在因子 2 上载荷小或可以忽略), 而变量 1 和 3 定义因子 2(在因子 2 上高载荷而在因子 1 上载荷小或可以忽略). 变量 4
